\section{Lepton mixing and a neutrino-scale interface}
\label{sec:pmns_neutrino_closure}

This section extends the bounded-complexity closure program to the lepton sector.
We record a minimal, auditable closure for the PMNS mixing angles and provide a corresponding matrix reconstruction in the PDG standard parameterization.
Majorana phases do not affect oscillation probabilities and are not constrained by this minimal closure, so we ignore them here \cite{PDGRPP2024}.

\subsection{PMNS angles as bounded-complexity amplitudes}
\label{subsec:pmns_mixing_depths}

We use the same golden resolution coordinate for dimensionless amplitudes as in \eqref{eq:r_mix_def}:
\[
r_{\mathrm{mix}}(x):=-\frac{\log x}{\log\varphi}\qquad(x>0).
\]

\paragraph{Motivation (normalization and depth).}
As in the CKM closure of Section~\ref{subsec:ckm_mixing_depths}, two generic finite-resolution mechanisms motivate the candidate family below:
(i) normalization in an effective $m$-dimensional subspace yields typical component magnitudes of order $1/\sqrt{m}$ (since squared magnitudes sum to $1$), and
(ii) discrete golden-branch depth costs yield exponentially small weights with base $\varphi$, whose square roots naturally produce amplitude factors $\varphi^{-k/2}$.
We treat the resulting PMNS hierarchies as a bounded-complexity closure problem rather than as free continuous inputs \cite{PDGRPP2024,Esteban2020NuFIT}.
For definiteness, the reference targets used by the accompanying scripts are taken as representative normal-ordering global-fit central values under the PDG convention \cite{PDGRPP2024,Esteban2020NuFIT}.

We define a simple candidate family for the three sines $(s_{12},s_{23},s_{13})$ of the PMNS mixing angles:
\[
s_{12}=\sqrt{\frac{p_{12}}{q_{12}}},\qquad
s_{23}=\sqrt{\frac{p_{23}}{q_{23}}},\qquad
s_{13}=\varphi^{-k_{13}/2},
\]
with a bounded-complexity box
\[
\begin{aligned}
&1\le q_{12},q_{23}\le B,\qquad 1\le p_{12}\le q_{12}-1,\qquad 1\le p_{23}\le q_{23}-1,\\
&\gcd(p_{ij},q_{ij})=1,\qquad 1\le k_{13}\le 2B,
\end{aligned}
\]
and we select the unique minimizer by lexicographic minimization of the maximum absolute log mismatch and then the sum mismatch (as in Definition~\ref{def:bounded_complexity_closure}).

\begin{proposition}[Bounded-complexity rigidity for PMNS mixing sines]\label{prop:pmns_mixing_rigidity}
Fix representative global-fit reference values for $(s_{12},s_{23},s_{13})$ under PDG conventions \cite{PDGRPP2024,Esteban2020NuFIT}.
For each $B\in\{2,\dots,20\}$, minimize the minimax log-mismatch objective over the candidate box described above, using the same deterministic tie-break rules as in Definition~\ref{def:bounded_complexity_closure}.
At $B=20$ the unique minimizer is
\[
(p_{12}/q_{12},p_{23}/q_{23},k_{13})=(4/13,6/11,8),
\]
and, within the tested range $B=2,\dots,20$, the minimizer stabilizes at $B=13$ and remains constant for all $B$ with $13\le B\le 20$ (Table~\ref{tab:pmns_mixing_rigidity}).
\end{proposition}
\begin{proof}
Finite exhaustive enumeration with deterministic tie-break rules; reproduced by the accompanying script.
\end{proof}

\begin{remark}[Candidate-domain size]\label{rem:pmns_domain_size}
Let $R(B)$ be the set of reduced fractions $p/q$ with $1\le q\le B$, $1\le p\le q-1$, and $\gcd(p,q)=1$.
Then $|R(B)|=\sum_{q=2}^{B}\varphi_E(q)$, where $\varphi_E$ is Euler's totient function \cite{HardyWright2008}.
The PMNS candidate box has size
\[
|\Theta(B)|=|R(B)|^2\cdot (2B),
\]
corresponding to independent choices for $(p_{12}/q_{12})$, $(p_{23}/q_{23})$, and $k_{13}\in\{1,\dots,2B\}$.
At $B=20$, one has $|R(20)|=127$ and therefore $|\Theta(20)|=127^2\cdot 40=645160$.
\end{remark}

\begin{table}[t]
\centering
\small
\setlength{\tabcolsep}{6pt}
\renewcommand{\arraystretch}{1.15}
\begin{tabular}{rllr}
\toprule
$B$ & minimizer & max abs.\ log mismatch & sum abs.\ log mismatch \\
\midrule
2 & $(p_{12}/q_{12},p_{23}/q_{23},k_{13})=(1/2,1/2,4)$ & 0.950499 & 1.237468 \\
3 & $(p_{12}/q_{12},p_{23}/q_{23},k_{13})=(1/3,1/2,6)$ & 0.469287 & 0.553524 \\
4 & $(p_{12}/q_{12},p_{23}/q_{23},k_{13})=(1/3,1/2,8)$ & 0.043089 & 0.096161 \\
5 & $(p_{12}/q_{12},p_{23}/q_{23},k_{13})=(1/3,1/2,8)$ & 0.043089 & 0.096161 \\
6 & $(p_{12}/q_{12},p_{23}/q_{23},k_{13})=(1/3,1/2,8)$ & 0.043089 & 0.096161 \\
7 & $(p_{12}/q_{12},p_{23}/q_{23},k_{13})=(2/7,4/7,8)$ & 0.035928 & 0.071529 \\
8 & $(p_{12}/q_{12},p_{23}/q_{23},k_{13})=(2/7,4/7,8)$ & 0.035928 & 0.071529 \\
9 & $(p_{12}/q_{12},p_{23}/q_{23},k_{13})=(2/7,5/9,8)$ & 0.035928 & 0.057444 \\
10 & $(p_{12}/q_{12},p_{23}/q_{23},k_{13})=(3/10,5/9,8)$ & 0.011925 & 0.033049 \\
11 & $(p_{12}/q_{12},p_{23}/q_{23},k_{13})=(3/10,6/11,8)$ & 0.011925 & 0.023874 \\
12 & $(p_{12}/q_{12},p_{23}/q_{23},k_{13})=(3/10,6/11,8)$ & 0.011925 & 0.023874 \\
13 & $(p_{12}/q_{12},p_{23}/q_{23},k_{13})=(4/13,6/11,8)$ & 0.011925 & 0.013468 \\
14 & $(p_{12}/q_{12},p_{23}/q_{23},k_{13})=(4/13,6/11,8)$ & 0.011925 & 0.013468 \\
15 & $(p_{12}/q_{12},p_{23}/q_{23},k_{13})=(4/13,6/11,8)$ & 0.011925 & 0.013468 \\
16 & $(p_{12}/q_{12},p_{23}/q_{23},k_{13})=(4/13,6/11,8)$ & 0.011925 & 0.013468 \\
17 & $(p_{12}/q_{12},p_{23}/q_{23},k_{13})=(4/13,6/11,8)$ & 0.011925 & 0.013468 \\
18 & $(p_{12}/q_{12},p_{23}/q_{23},k_{13})=(4/13,6/11,8)$ & 0.011925 & 0.013468 \\
19 & $(p_{12}/q_{12},p_{23}/q_{23},k_{13})=(4/13,6/11,8)$ & 0.011925 & 0.013468 \\
20 & $(p_{12}/q_{12},p_{23}/q_{23},k_{13})=(4/13,6/11,8)$ & 0.011925 & 0.013468 \\
\bottomrule%
\end{tabular}
\caption{Bounded-complexity minimax search for PMNS mixing sines in the candidate family described in the text. Rows are generated by the reproducible script \texttt{scripts/exp\_pmns\_mixing\_depth\_rigidity.py}. Reference targets use representative global-fit central values under PDG conventions \cite{PDGRPP2024,Esteban2020NuFIT}.}
\label{tab:pmns_mixing_rigidity}
\end{table}

\begin{table}[t]
\centering
\small
\setlength{\tabcolsep}{6pt}
\renewcommand{\arraystretch}{1.15}
\begin{tabular}{lrrrrr}
\toprule
parameter & reference & predicted form & predicted value & $r_{\mathrm{mix}}$ & $\log(\mathrm{pred}/\mathrm{ref})$ \\
\midrule
$s_{12}$ & 0.554076 & $\sqrt{4/13}$ & 0.554700196 & 1.225 & 0.001126 \\
$s_{23}$ & 0.738241 & $\sqrt{6/11}$ & 0.738548946 & 0.630 & 0.000417 \\
$s_{13}$ & 0.147648 & $\varphi^{-8/2}$ & 0.145898034 & 4.000 & -0.011925 \\
\bottomrule%
\end{tabular}
\caption{PMNS mixing-angle closure in the golden resolution coordinate at $(m,n)=(6,3)$, using the bounded-complexity minimizer at $B=20$. Rows are generated by \texttt{scripts/exp\_pmns\_mixing\_depth\_rigidity.py}.}
\label{tab:pmns_mixing}
\end{table}

\paragraph{Audit context.}
Candidate-domain size, best/second-best mismatch context, and distribution quantiles for the PMNS-sine closure (and the bounded-denominator $\delta$ closure below) are recorded in Appendix~\ref{app:generated_tables} (Tables~\ref{tab:audit_closure_metrics}--\ref{tab:audit_closure_quantiles}); uncertainty-robustness stress tests are recorded in Table~\ref{tab:audit_uncertainty_robustness}, and counterfactual baseline comparisons in Table~\ref{tab:audit_counterfactual}.

\subsection{Matrix reconstruction and a discrete CP-phase closure}
\label{subsec:pmns_matrix_closure}

To reconstruct a full $3\times 3$ mixing matrix, we use the PDG standard parameterization (three angles plus one Dirac CP phase) \cite{PDGRPP2024}.
We compare a single reference reconstruction (from representative global-fit central values) to a closed prediction obtained from the bounded-complexity angle minimizer together with a discrete protocol-level closure for the Dirac phase $\delta$ \cite{Esteban2020NuFIT}.
Because the PMNS \emph{magnitudes} depend on $\cos\delta$ while the CP-odd invariant depends on $\sin\delta$, we use a bounded-denominator rational-angle candidate family
\[
\delta=\frac{k\pi}{q},\qquad 1\le q\le Q,
\qquad 1\le k\le 2q-1,\qquad \gcd(k,q)=1,
\]
and we select a unique minimizer by a CP-odd anchor rule:
first require $\mathrm{sgn}(J_\ell)$ to match the sign induced by the reference reconstruction, then minimize the mismatch $|\log(|J_\ell|/|J_{\ell,\mathrm{ref}}|)|$, and finally break ties by bounded-complexity order $(q,k)$.
To remove the $\sin\delta$ sign ambiguity \emph{and} fix the correct quadrant (since magnitudes depend on $\cos\delta$), we impose a second anchor:
we require $\mathrm{sgn}(\cos\delta)$ to match the sign induced by the same reference reconstruction.
Table~\ref{tab:pmns_delta_sweep} records the resulting bounded-denominator sweep $Q=1,\dots,12$ and the stabilized minimizer.

\begin{table}[t]
\centering
\small
\setlength{\tabcolsep}{10pt}
\renewcommand{\arraystretch}{1.15}
\resizebox{\textwidth}{!}{%
\begin{tabular}{rllrrrr}
\toprule
$Q$ & minimizer & $\delta$ [deg] & $\mathrm{sgn}(J_\ell)$ match & $\mathrm{sgn}(\cos\delta)$ match & $J_\ell$ (closed) & $|\log(|J_\ell|/|J_{\ell,\mathrm{ref}}|)|$ \\
\midrule
90.0 & FLIP & +0.0336565 & 1.366 \\
\textbf{270.0} & OK & -0.0336565 & 1.366 \\
\bottomrule%
\end{tabular}
}%
\caption{Bounded-denominator closure for the PMNS Dirac phase $\delta$ using a CP-odd anchor and a deterministic tie-break by bounded complexity. Rows are generated by \texttt{scripts/exp\_pmns\_matrix\_closure.py}.}
\label{tab:pmns_delta_sweep}
\end{table}

\begin{table}[t]
\centering
\small
\setlength{\tabcolsep}{6pt}
\renewcommand{\arraystretch}{1.15}
\begin{tabular}{lrrr}
\toprule
parameter & reference value & closed value & mismatch \\
\midrule
$s_{12}$ & 0.554075807 & 0.577350269 & 0.041148 \\
$s_{23}$ & 0.738241153 & 0.707106781 & -0.043089 \\
$s_{13}$ & 0.147648231 & 0.145898034 & -0.011925 \\
$\delta$ [deg] & 195.000000 & 270.000000 & 75.000000 \\
$J_\ell$ & -0.00858602 & -0.0336565 & 1.366071 \\
\bottomrule%
\end{tabular}
\caption{PMNS angle/phase extraction and the induced leptonic Jarlskog invariant $J_\ell$ under the chosen conventions. Rows are generated by \texttt{scripts/exp\_pmns\_matrix\_closure.py}.}
\label{tab:pmns_angles}
\end{table}

\begin{table}[t]
\centering
\small
\setlength{\tabcolsep}{6pt}
\renewcommand{\arraystretch}{1.15}
\begin{tabular}{lrrr}
\toprule
element & reference $|U_{ij}|$ & closed $|U_{ij}|$ & $\log(\mathrm{closed}/\mathrm{ref})$ \\
\midrule
$|U_{e1}|$ & 0.823342335 & 0.807759768 & -0.019107 \\
$|U_{e2}|$ & 0.548003102 & 0.571172409 & 0.041410 \\
$|U_{e3}|$ & 0.147648231 & 0.145898034 & -0.011925 \\
$|U_{\mu1}|$ & 0.287059885 & 0.416847788 & 0.373030 \\
$|U_{\mu2}|$ & 0.620062595 & 0.580414541 & -0.066078 \\
$|U_{\mu3}|$ & 0.730149985 & 0.699540479 & -0.042826 \\
$|U_{\tau1}|$ & 0.489595774 & 0.416847788 & -0.160859 \\
$|U_{\tau2}|$ & 0.561440093 & 0.580414541 & 0.033237 \\
$|U_{\tau3}|$ & 0.667143913 & 0.699540479 & 0.047418 \\
\bottomrule%
\end{tabular}
\caption{PMNS magnitude table induced by the bounded-complexity angle closure and the discrete phase choice. Rows are generated by \texttt{scripts/exp\_pmns\_matrix\_closure.py}.}
\label{tab:pmns_matrix}
\end{table}

\begin{table}[t]
\centering
\small
\setlength{\tabcolsep}{10pt}
\renewcommand{\arraystretch}{1.15}
\begin{tabular}{lrr}
\toprule
unitarity check & reference value & closed value \\
\midrule
row 1 & +0.000e+00 & +2.220e-16 \\
row 2 & -1.110e-16 & +2.220e-16 \\
row 3 & +0.000e+00 & +2.220e-16 \\
col 1 & +0.000e+00 & +2.220e-16 \\
col 2 & +0.000e+00 & +2.220e-16 \\
col 3 & -1.110e-16 & +0.000e+00 \\
\bottomrule%
\end{tabular}
\caption{Row/column unitarity diagnostics: $\sum_j |U_{ij}|^2-1$ (rows) and $\sum_i |U_{ij}|^2-1$ (columns) for the reconstructed and closed PMNS matrices.}
\label{tab:pmns_unitarity}
\end{table}

\subsection{A minimal neutrino mass-scale interface}
\label{subsec:neutrino_mass_interface}

At the minimal balanced resolution $(m,n)=(6,3)$, the mass-spectrum closure of Section~\ref{sec:mass_spectrum_closure} anchors scheme-stable charged-lepton scales and treats neutrino absolute masses as an interface input.
To express neutrino scales in the same resolution language, we record a deterministic nearest-integer depth assignment for representative minimal-mass normal/inverted orderings inferred from oscillation splittings.
We use representative central values for $(\Delta m_{21}^2,|\Delta m_{31}^2|)$ from standard global fits \cite{PDGRPP2024,Esteban2020NuFIT}.

\begin{table}[t]
\centering
\small
\setlength{\tabcolsep}{6pt}
\renewcommand{\arraystretch}{1.15}
\begin{tabular}{llrrrrr}
\toprule
ordering & eigenstate & reference $\mu$ [GeV] & $r(\mu)$ & nearest $\widehat r$ & $\Delta r$ & $\mu/\mu_{\mathrm{pred}}$ \\
\midrule
NO & $m_1$ & $0$ & $-$ & $-$ & $-$ & $-$ \\
NO & $m_2$ & $8.61394e-12$ & -37.195 & -37 & -0.195 & $0.910594$ \\
NO & $m_3$ & $5.01697e-11$ & -33.533 & -34 & 0.467 & $1.25199$ \\
IO & $m_1$ & $4.998e-11$ & -33.541 & -34 & 0.459 & $1.24726$ \\
IO & $m_2$ & $5.07169e-11$ & -33.510 & -34 & 0.490 & $1.26565$ \\
IO & $m_3$ & $0$ & $-$ & $-$ & $-$ & $-$ \\
\bottomrule%
\end{tabular}
\caption{A minimal neutrino mass-scale interface in the resolution coordinate. For each reference mass $\mu$ we compute $r(\mu)=\log(\mu/m_e)/\log\varphi$ and select the nearest integer depth $\widehat r$, yielding a depth mismatch $\Delta r=r-\widehat r$ and the implied multiplicative matching factor. Rows are generated by \texttt{scripts/exp\_neutrino\_mass\_interface.py}.}
\label{tab:neutrino_mass_interface}
\end{table}


